\section*{第 \ref{ch:b3:conservation-biology} 章——保护生物学}

\begin{solution}{exo:b3:conservation-biology:1}
\omterm{def:b3:conservation-biology:small}{种群统计随机性}：少数个体的出生、死亡与性别
所带来的偶然波动——如数量为 51 只的卡卡坡鹦鹉，其中雄性多于
雌性。\omterm{def:b3:conservation-biology:small}{环境随机性}：某个不利年份同时打击所有个体——
如干旱、疾病或严冬，袭击一个仅有几百个体的
种群。\omterm{def:b3:conservation-biology:small}{阿利效应}：低密度时人均增长率下降——旅鸽的
群体在数量稀疏后便无法繁殖，其残存的
捕食者也不再
有大量猎物可供捕食。
\end{solution}

\begin{solution}{exo:b3:conservation-biology:2}
理想种群的大小——规模恒定、随机交配、
两性数量相等、家系大小服从泊松分布——其\omterm{thm:b3:conservation-biology:ne}{杂合度}丧失的速率
与实际种群相同。繁殖雄性与雌性的数量不等；
种群规模的波动使最小的世代权重更大；
家系大小的方差高于泊松分布，少数个体
产下大部分后代（另有非随机交配与
世代重叠）。
\end{solution}

\begin{solution}{exo:b3:conservation-biology:3}
$S = cA^{z}$。保留比例 $= (A'/A)^{z}$：$0.5^{0.25} = 0.84$；
$0.1^{0.25} = 0.56$；$0.01^{0.25} = 0.32$——面积缩至百分之一，
最终仍能保有约三分之一的物种。
\end{solution}

\begin{solution}{exo:b3:conservation-biology:4}
$N_{e} \ge 50$ 可使每代近交增量低于 $1/100$，
即动物育种者可容忍的近交衰退速率，
它可防止短期内的适合度下降。$N_{e} \ge
500$ 则使漂变造成的变异丧失与突变
所补充的变异相互平衡，使种群保有适应所需的
数量变异：它保障的是长期适应。实际种群的
普查数量需为上述数值的五至十倍。
\end{solution}

\begin{solution}{exo:b3:conservation-biology:5}
$N_{e} = 4\times 8\times 120/128 = 30$；\omterm{thm:b3:conservation-biology:ne}{杂合度}每代下降 $1/60 =
\qty{1.7}{\%}$；损失四分之一时，$(1 - 1/60)^{t} =
0.75$，$t = \ln 0.75/\ln(59/60) = 17$ 代。
\end{solution}

\begin{solution}{exo:b3:conservation-biology:6}
$1/N_{e} = \tfrac{1}{5}(4/1000 + 1/20) = 0.0108$，$N_{e} = 93$：
规模为 20 的一代使五代的行为如同规模为 93 的种群。
\omterm{thm:b3:conservation-biology:ne}{杂合度}保留率为 $(1 - 1/186)^{5} = 0.973$，相较于 $(1 -
1/2000)^{5} = 0.9975$，即恒定规模为 1000 的种群的保留率。
\end{solution}

\begin{solution}{exo:b3:conservation-biology:7}
$F = 1 - (1 - 1/20)^{5} = 0.23$；适合度下降 $5\times 2.3 = \qty{11}{\%}$。
得克萨斯雌性与佛罗里达雄性的每个后代的 $F = 0$；若
八只新迁入者占繁殖雌性的五分之二，则下一代中的五分之二
为远交个体，平均 $F$ 降至约
$0.6\times 0.23 = 0.14$，一代之内——这正是
恢复过程所显示的结果。
\end{solution}

\begin{solution}{exo:b3:conservation-biology:8}
$r = \ln(100/31)/8 = \qty{0.15}{yr^{-1}}$。若延续至 2020 年：
$100\,\mathrm{e}^{0.15\times 17} \approx 1200$。公园内约维持
一百只：领地已被占满，麋鹿数量下降，群体相互杀害
成员，疥螨病与犬瘟热随之蔓延——密度制约
在空旷分布区被占满后开始显现。
\end{solution}

\begin{solution}{exo:b3:conservation-biology:9}
(a) 十年内下降 \qty{60}{\%} 超过 \qty{50}{\%}：濒危。
(b) 成熟个体数 $180 < 250$：极危。(c) 分布范围
\qty{3000}{km^2}，低于 \qty{5000}{}：濒危。(d) \num{30000}
个体，下降 \qty{10}{\%}：两项阈值均未达到——不
受威胁（至多为近危，需持续关注）。
\end{solution}

\begin{solution}{exo:b3:conservation-biology:10}
全部 $N$ 对繁殖失败的概率为 $(1/2)^{N}$：$0.25$，$0.031$，
$0.001$，$10^{-6}$。种群统计风险随 $N$ 呈指数下降，
在几十个个体以上可忽略不计；超过此规模后，残留的风险
为环境风险，其下降并不那么快，因为它
同时打击所有个体。
\end{solution}

\begin{solution}{exo:b3:conservation-biology:11}
占十分之一的一个片段，其保留原有物种的比例为 $0.1^{0.25} = 0.56$。
十个片段所保留的比例介于 $0.56$（若各片段物种相同）与
理论上可超过整体（若各片段各保有不同物种——
超出区域物种库则不可能）。它们落在何处取决于
片段间的周转：不同片段中的不同生境
会使总量上升；需要大面积、内部生境的物种，或
因相同\omterm{prop:b3:conservation-biology:area}{边缘效应}而在各片段中消失的物种，会使其
下降；且各片段的灭绝债务分别偿还，因此
长期总量通常低于整体的物种数。
\end{solution}

\begin{solution}{exo:b3:conservation-biology:12}
漂变使\omterm{thm:b3:conservation-biology:ne}{杂合度}以每代 $1/2N_{e}$ 的速率消失；
迁入者以取自别处的等位基因替换基因库中
比例为 $m$ 的部分。若 $mN_{e} \approx 1$，$m \approx 1/N_{e}$，为丧失速率的两倍，
因此变异补充的速度快于其衰减，使
种群的等位基因频率与源种群保持接近（$F_{ST} = 1/(1
+ 4N_{e}m) \approx 0.2$）。代价是：局部有利的等位基因
每代被 $m$ 所稀释，因此局部适应仅存在于
选择系数超过约 $1/N_{e}$ 的性状；弱选择的
局部差异则被淹没。
\end{solution}

\begin{solution}{pb:b3:conservation-biology:1}
\textbf{1.} $N_{e} = 4\times 20\times 40/60 = 53$，而普查数为 60。
\textbf{2.} $1/N_{e} = \tfrac{1}{5}(1/500 + 1/100 + 1/60 + 1/51 + 1/60) =
0.0130$，$N_{e} = 77$；规模为 51 与 60 的瓶颈世代
起主导作用，500 的影响几乎可忽略。
\textbf{3.} $F = 1 - (1 - 1/154)^{5} = 0.032$；线性的结果为 $5/154 = 0.032$。
\textbf{4.} $6\times 0.32 = \qty{1.9}{\%}$.
\textbf{5.} $(1 - 1/154)^{5} = 0.968$：\qty{96.8}{\%} 保留。
\textbf{6.} 在 $N_{e} = 53$ 下经历十代：增量 $1 - (1 -
1/107)^{10} = 0.090$，总 $F = 1 - 0.968\times 0.910 = 0.12$；适合度
下降 $6\times 1.2 = \qty{7}{\%}$；历时一个世纪。
\textbf{7.} 本地–本地配对占交配的 $(60/70)^{2} = 0.73$，
其后代保持 $F \approx 0.03$（另加少量增量）；其余
配对的后代 $F = 0$：平均 $F \approx 0.73\times 0.035 = 0.026$，
一代之内下降约四分之一。
\textbf{8.} $F$ 是个体内的同源一致性，任何
远交后代都会使其归零；\omterm{thm:b3:conservation-biology:ne}{杂合度}是
等位基因频率的性质，迁入者的等位基因只有在
被传递下去时才会扩散。种群仍然很小时，漂变
会在数十代内再次将其清除，因此拯救必须重复进行，或
使种群扩大。
\textbf{9.} 两性数量相等时 $N_{e} = N$，成体 500 个；以 70 为起点，$r =
0.05$：$\ln(500/70)/0.05 = 39$ 年。
\textbf{10.} 理想种群的家系大小服从泊松分布，方差
等于均值 2；$N_{e} = 4N/(2 + V_{k})$，因此使每一对
贡献相等（$V_{k} = 0$）则得 $N_{e} = 2N$：没有任何谱系
会因巢中的运气而消失。
\textbf{11.} 采集并保存其精液，用于人工授精，
并在数年间用于多只无亲缘关系的雌性——代价是
其等位基因可能主导一个小的基因库，因此其份额必须
限制在 $1/N_{e}$ 附近；或将组织冷冻保存，以备日后克隆或
配子的产生，代价是延迟与技术风险。若不作为，
该谱系的等位基因将永久丢失。
\textbf{12.} 约 1.6 亿美元用于约 430 只鸟，每只约
\num{370000} 美元。按单只鹦鹉计算昂贵，相对一个物种则便宜——
且不及几公里高速公路的造价，
这正是预算实际采用的比较。
\textbf{13.} $(300/2000)^{0.25} = 0.62$：约保留 180 种中的 112 种，
灭绝债务为 68。
\textbf{14.} 每个片段含 $(100/2000)^{0.25}\times 180 = 85$ 个物种。
若各片段完全不同：至多 255，超出 180 的区域物种库则不可能；
若完全相同：85。决定因素在于各片段生境
的差异程度，以及 \qty{100}{km^2} 内能维持多少物种；现实中
介于 85 与 112 之间。
\textbf{15.} $(320/2000)^{0.25}\times 180 = 114$，对比三个孤立相同片段的
85：\omterm{prop:b3:conservation-biology:area}{廊道}使再定殖与基因流
令三个小种群表现得如同一个更大的种群。
\textbf{16.} 两性相等时 $N_{e} = 50$ 对应 25 对，需 \qty{1250}{km^2}。
整个 \qty{300}{km^2} 可容纳 6 对，$N_{e} \approx 12$；一个
片段则为 2 对。两者都不够：捕食者无论
物种数多少都会消失，保护其所需的面积即可保护所有
更小的物种——即伞护种。
\textbf{17.} $\ln(200/10)/0.03 = 100$ 年：债务需经过一个
世纪方能偿清，远在森林被砍伐之后。
\textbf{18.} \omterm{prop:b3:conservation-biology:area}{边缘效应}使每个片段丧失内部生境，因此
可用面积小于图上所示面积，损失
大于定律所预测的程度；次生生长的基质或树篱
若某些物种能加以利用，则增加有效面积，使损失变小。
\textbf{19.} $r = \ln(100/31)/8 = \qty{0.146}{yr^{-1}}$；倍增时间
$\ln 2/0.146 = 4.7$ 年。
\textbf{20.} $N(8) = 120/[1 + (120/31 - 1)\mathrm{e}^{-1.17}] = 120/(1 +
2.87\times 0.31) = 63$。观测值 100 即指数模型的数值本身——
$r$ 正是据这两个点拟合的——与逻辑斯谛模型的
63 相比：头几年狼群未遇到密度制约，
领地空闲，驼鹿充足。
\textbf{21.} $\ln(4000/\num{17000})/20 = -0.072$：每年 \qty{7}{\%}。一百只
狼每年从鹿群中捕食 \num{2000} 头，该鹿群平均
\num{8000} 头，即每年四分之一，单凭此已足以造成下降，
尽管园区外的狩猎、干旱与灰熊对幼崽的捕食
也分担了这一下降。
\textbf{22.} 成熟个体仅六十只，低于 250：极危。
捕食者 $N_{e} \approx 12$（仅数十只）：同样
列为极危。
\textbf{23.} 由 $3\times 10^{9}$ 降至 $1$ 历时 44 年：$r = -\ln(3\times
10^{9})/44 = -0.50$ 每年，每十七个月减半。没有任何
捕获能每年移除数十亿之半；下降起初较慢，
末期则快得多，因为群居繁殖者密度低于阈值时便停止繁殖，
而其捕食者不再被淹没——
\omterm{def:b3:conservation-biology:small}{阿利效应}使衰退变成了崩溃。
\textbf{24.} 1850 年时种群数量以十亿计，其实测增长率为正，
方差小，而\omterm{met:b3:conservation-biology:pva}{种群生存力分析}只外推所给定的动态：
它会把风险估为零。
它会漏掉外部驱动因素——铁路与电报使猎人得以追踪
并清空每一处栖息地，还有森林砍伐——以及
群居繁殖者的阿利阈值，两者都不在
丰富物种的种群统计之中；\omterm{met:b3:conservation-biology:pva}{种群生存力分析}只是外推，灭绝的原因
往往是新出现的。
\textbf{25.} 今日 $N_{e} \approx 53$；再经十代后 $F \approx 0.12$，
每代 60 只，适合度降低 \qty{7}{\%}；缩小后的森林保有
约 \qty{62}{\%} 的物种。
\end{solution}
